% Transcription — fonds Alexandre Grothendieck, shelfmark 32, batch 3 (pages 41-47).
% Established from the facsimile archives/batches/32/batch-03.pdf (a mirror of
% https://grothendieck.umontpellier.fr/32.pdf), per the skill
% transcribe-grothendieck.
%
% Pass: Opus 5.5 (claude-opus-5-5), 2026-09-24 — first pass, unchecked against the pages by a human.
%
% The batch is the tail of a manuscript section 3.1 on the functor Rf^! for a
% compactifiable morphism of schemes (numbered paragraphs 3.1.14 to 3.1.19),
% in blue ink on squared paper, renumbered in places in green and black ink,
% and annotated throughout in pencil by a reader asking for references,
% querying notation and proposing a new plan for 3.1 (page 47). The two
% writings — ink body, pencil annotations — are kept apart: the pencil goes
% into \marginal{} (or \add{} where it is an insertion into the running
% text), and the file does not say whose hand either is.
%
%   Page 41 — 3.1.14: Rf^! computed by K(U) = f̄_* τ_{≤d} C^* A_U = f_!^•(A_U);
%     for nA = 0, Rf_! A_U ≃ A ⊗^L Rf_!(Z/n)_U; flatness of f_!^•((Z/n)_U)
%     (reduction to n = ℓ^k), so f_!^•(Z/n)_U ⊗ A → f_!^•(A_U) is a
%     quasi-isomorphism.
%   Page 42 — the canonical flat resolution L_•(F); Rf^! given by
%     U ↦ L_•(f_!^•(Z/n)_U); the ring sheaf A plays only a dummy role.
%   Page 43 — a narrow slip: he has not checked that this construction agrees
%     with that of 3.1.9.
%   Page 44 — 3.1.15: Rf^! by the procedure of 3.1.13, K(W) for
%     W = (U, V, φ) ∈ Ob X_f; a pencil diagram in the margin (redrawn).
%   Page 45 — a single struck line, « Proposition 3.1.12. Le morphis ».
%   Page 46 — 3.1.16: the adjunction isomorphism at the level of complexes and
%     its localised form (3.1.16.2); 3.1.17, Exemple: stalks of
%     H^{-i}(Rf^! Z/nZ) at a closed point of X over an algebraically closed
%     field, and the bidualité theorem (SGA 5 I, reference left blank).
%   Page 47 — end of 3.1.17 (pro-objects); Proposition 3.1.18 (amplitude and
%     injective dimension of Rf^!L); 3.1.19 (quasi-finite f, immersions,
%     R^i f^! L = H^i_X(L)); pencil: « Le plan de 3.1 semble flou après
%     3.1.4 » with a three-part plan, and a pencil computation of f^! for
%     smooth f (base change + relative purity), with a small diagram.
%
% Nothing skipped: every page of the batch carries mathematics. Page 45 is
% given for its one struck line; page 43 is a slip of two lines.
%
% The dating is the inventory's own, for the whole shelfmark.
\documentclass[11pt,a4paper]{article}
% Le préambule commun à toutes les transcriptions du fonds Grothendieck.
%
% Il tient en une idée : **l'incertitude se compose, elle ne se résout pas.**
% Une transcription qui lisse un mot douteux a détruit la seule chose pour
% laquelle on la faisait — un lecteur doit pouvoir savoir, à chaque endroit,
% ce qui était sur le feuillet et ce qui a été ajouté. Les six macros qui
% suivent sont donc l'appareil critique, et non de la décoration.
%
% Les mêmes commandes sont reconnues par `scripts/render.mjs`, qui produit la
% vue de lecture du site. Une macro ajoutée ici doit l'être aussi là-bas,
% faute de quoi le rendu échouera bruyamment — ce qui est voulu : un rendu
% silencieusement faux est pire qu'un rendu absent.

\ProvidesPackage{grothendieck}[2026/08/08 Transcriptions du fonds Grothendieck]

\usepackage{amsmath,amssymb,amsthm}
\usepackage{mathrsfs}
\usepackage{stmaryrd}
% Les diagrammes commutatifs sont partout dans ce fonds : c'est la langue
% même de la géométrie algébrique de Grothendieck, et une transcription qui
% les rendrait en prose ne transcrirait rien.
\usepackage{tikz-cd}
% Français par défaut — c'est la langue des feuillets — mais l'anglais est
% chargé aussi : la traduction et le résumé sont des documents anglais, et la
% typographie française (espaces avant les deux-points) y serait une faute.
% Ces documents basculent par \selectlanguage{english} en tête de corps.
\usepackage[english,french]{babel}
\usepackage{fontspec}
\usepackage{geometry}
\geometry{a4paper,margin=25mm,marginparwidth=28mm}
\usepackage{marginnote}
\usepackage{xcolor}
% ulem, not soul. `soul`'s \st and \ul are built on a character-by-character
% scan that XeLaTeX's font machinery breaks: the document fails with an
% undefined control sequence rather than degrading. ulem does the same two jobs
% and is Unicode-engine safe. `normalem` stops it hijacking \emph, which the
% transcriptions use for Grothendieck's own emphasis.
\usepackage[normalem]{ulem}
\usepackage{hyperref}
\hypersetup{colorlinks=true,linkcolor=black,urlcolor=black,pdfborder={0 0 0}}

% --- Métadonnées du lot -----------------------------------------------------
% Reprises telles quelles par le script de rendu, qui les affiche en tête de
% la vue de lecture. Elles fixent ce que le fichier prétend couvrir : sans
% elles, une transcription flotte sans point d'attache dans un fonds de seize
% mille feuillets.

\newcommand{\folder}[1]{\gdef\grothFolder{#1}}
\newcommand{\batch}[1]{\gdef\grothBatch{#1}}
\newcommand{\leaves}[2]{\gdef\grothFirst{#1}\gdef\grothLast{#2}}
\newcommand{\foldertitle}[1]{\gdef\grothFoldertitle{#1}}
\gdef\grothFolder{?}\gdef\grothBatch{?}\gdef\grothFirst{?}\gdef\grothLast{?}\gdef\grothFoldertitle{}

% --- Le résumé --------------------------------------------------------------
%
% Il ouvre la lecture modernisée, sous le titre « Résumé », et il est composé
% en petit. Ce n'est pas un ornement : c'est le seul passage du document qui
% ne soit pas des mathématiques, et un lecteur doit pouvoir le reconnaître
% avant de l'avoir lu — puis le sauter s'il connaît déjà le sujet, ou s'y
% arrêter s'il ne le connaît pas. Le filet qui le suit marque la reprise du
% corps.

\newenvironment{resume}
  {\small\setlength{\parskip}{0.45em}}
  {\par\medskip\hrule\medskip}

% --- Mots-clés ---------------------------------------------------------------
%
% En anglais, en clôture du résumé de la lecture modernisée : le vocabulaire
% moderne sous lequel chercher ce que les feuillets construisent. Le manifeste
% du site (`npm run manifest`) les extrait pour étiqueter la cote entière —
% cette ligne est la source unique des tags, il n'y a pas de fichier à côté.
\newcommand{\keywords}[1]{\par\smallskip\noindent
  {\footnotesize\textsc{Keywords} --- \textit{#1}}\par}

% --- L'appareil critique ----------------------------------------------------

% Le numéro de feuillet, en marge. C'est l'ancre qui permet de régler un
% désaccord sur une formule en regardant *un* feuillet plutôt qu'une chemise
% de six cents. Le site s'en sert aussi pour tourner les pages du fac-similé
% quand on fait défiler la transcription.
\newcommand{\leaf}[1]{%
  \marginnote{\footnotesize\color{gray!70}\textsf{#1}}%
  \label{leaf:#1}\ignorespaces}

% Illisible : on ne devine pas. Un mot inventé qui se lit comme les autres est
% le pire résultat possible.
\newcommand{\ill}{\textcolor{red!60!black}{[\,\ldots\,]}}

% Lecture douteuse : le mot est donné, et signalé comme incertain.
\newcommand{\uncertain}[1]{\uline{#1}}

% Ajout de l'éditeur — une lettre manquante, un mot sous-entendu.
\newcommand{\add}[1]{\textcolor{blue!50!black}{[#1]}}

% Biffé par Grothendieck lui-même. Ce qu'il a rayé fait partie du document :
% une rature dit souvent où la pensée a changé de direction.
\newcommand{\struck}[1]{\sout{#1}}

% Note du transcripteur, sur l'état matériel du feuillet.
\newcommand{\note}[1]{\par\smallskip{\small\color{gray!60!black}\textit{[#1]}}\par\smallskip}

% Note marginale de Grothendieck, distincte du corps du texte.
\newcommand{\marginal}[1]{\par\smallskip\noindent\hspace{1em}%
  {\small\color{gray!60!black}#1}\par\smallskip}

% --- Titre ------------------------------------------------------------------

\AtBeginDocument{%
  \begin{center}
    {\large\bfseries Fonds Alexandre Grothendieck --- cote n\textsuperscript{o}~\grothFolder}\\[2pt]
    {\itshape\grothFoldertitle}\\[4pt]
    {\small feuillets \grothFirst--\grothLast\ (lot \grothBatch)}\\[2pt]
    {\footnotesize\color{gray!60!black}Transcription établie d'après le fac-similé numérisé
     par l'Université de Montpellier.\\
     Les lectures incertaines sont soulignées, les illisibles notées [\,\ldots\,],
     les ajouts entre crochets.}
  \end{center}
  \vspace{1em}}

\endinput


\folder{32}
\batch{3}
\pages{41}{47}
\dating{[vers 1963-1977]}
\watermark{Édition de démonstration}
\foldertitle{Corrections (références, rajouts) [dont SGA 4] : tapuscrit (s.d.), notes manuscrites (s.d.).}

\begin{document}

\page{41}

\note{Le texte est à l'encre bleue ; certains numéros de paragraphe et de
renvoi y sont repassés ou corrigés à l'encre verte ou noire, et l'on donne le
numéro qui se lit en dernier. Des annotations au crayon, d'une écriture
distincte du corps, entourent des passages, posent des questions et
proposent des insertions : on les donne en \emph{marginal} quand elles sont
en marge, en \emph{ajout} quand elles s'insèrent dans le texte, en le
signalant. Les nombreux soulignements à l'encre bleue (mots d'un énoncé,
références) ne sont pas reproduits. Dans ce lot, $f_!^{\bullet}$ désigne le
foncteur au niveau des complexes (un point en exposant sur la page) et
$f_!^{i}$ ses composantes ; la page distingue les deux, et on les distingue
comme elle.}

\marginal{À quoi sert cette section 3.1.14, ? qu'y fait-on ?}
\note{au crayon, en haut de la marge de gauche.}

3.1.14. Appliquons les considérations qui précèdent à un morphisme
compactifiable de schémas $f : X \to S$.

D'après les démonstrations de 3.1.\struck{\ill{}}3 et 3.1.4, si
$f = \bar f j$ est une compactification de $f$ et si les fibres de $f$ sont
de dimension $\leqslant d$, le foncteur $Rf^!$ se construit par le procédé de
3.1.10, pour
\[
K(U) = \bar f_{*}\, \tau^{*}_{\leqslant d}\, C_{\ell}^{*} \mathcal{A}_U
= f_!^{\bullet}(\mathcal{A}_U) .
\]
\add{où $A_U$ …}
\note{« où $A_U$ … » est écrit au crayon sous la formule, et laissé en
suspens. Le « 4 » de 3.1.14 et le « 1 » de 3.1.10 sont repassés à l'encre
noire ; dans 3.1.3, un chiffre est noirci avant le « 3 ». L'astérisque
au-dessus de $\tau$ est appuyé, et $\tau^{*}_{\leqslant d}$ est entouré au
crayon.}
\marginal{3.1.10}
\marginal{$\tau''_{\leqslant 2d}$ ?}
\marginal{dire qui est $A_U$}
\note{trois annotations au crayon, chacune entourée, dans la marge de gauche
en regard de la formule.}

Le corollaire 3.1.12 va permettre de donner d'autres procédés de calcul du
foncteur $Rf^!$. Supposons que $n\mathcal{A} = 0$. On a alors
(XVII 5.2)
\[
Rf_!\, \mathcal{A}_U \simeq \mathcal{A} \overset{\mathbb{L}}{\otimes} Rf_!(\mathbb{Z}/n)_U .
\]
\note{le « 2 » de 3.1.12 est à l'encre verte, sur un autre chiffre. Au
crayon, un signe d'insertion suit « 5.2 », et un signe semblable est porté
dans la marge, sans texte.}

En fait, $f_!^{\bullet}((\mathbb{Z}/n)_U)$ est un complexe de faisceaux plats :
il suffit de le voir pour $n$ \add{de la forme} $\ell^k$ \struck{primaire} ;
\struck{\ill{}} la platitude \struck{\ill{}} d'un faisceau $\underline{F}$
\struck{\ill{}} de $\mathbb{Z}/\ell^k$-modules signifie que pour $i < k$,
\[
\underline{F}/\ell \underline{F} \xrightarrow{\ \sim\ } \ell^{i}\underline{F}/\ell^{i+1}\underline{F} ,
\]
et ceci résulte pour les faisceaux $f_!^{i}(\mathbb{Z}/n)_U$ de l'exactitude
de $f_!^{i}$ et de la platitude \struck{\ill{}} $(\mathbb{Z}/n)_U$.
\note{« de la forme » est ajouté au crayon, qui biffe aussi « primaire » et
souligne « plats ». Le mot biffé devant $(\mathbb{Z}/n)_U$ est une notation
en $f_!$ recouverte d'encre.}
\struck{On a donc un quasi-isomorphisme} La flèche \add{canonique}
\[
\struck{f_!^{\bullet}(\mathbb{Z}/n)_U \xrightarrow{\ \sim\ }}
\]
\[
f_!^{\bullet}(\mathbb{Z}/n)_U \otimes \mathcal{A} \xrightarrow{\ \sim\ } f_!^{\bullet}(\mathcal{A}_U)
\]
est donc un quasi-isomorphisme et on peut aussi bien décrire (3.1.12)
\note{« canonique » est au crayon. Le « 2 » de (3.1.12) est à l'encre verte.}

\page{42}

le foncteur $Rf^!$ comme fourni par la \struck{\ill{}} donnée 3.1.10
\[
U \longmapsto f_!^{\bullet}((\mathbb{Z}/n)_U)
\]
\note{le « 0 » de 3.1.10 est à l'encre verte ; la formule est entourée au
crayon et reliée à la note de marge qui suit.}
\marginal{$\otimes\mathcal{A}$ ? $\mathbb{Z}/n$ \uncertain{sinon}, expliciter
le \uncertain{nouveau} $\mathcal{A}$, $\mathcal{B}$}

Rappelons que la résolution plate canonique \add{$L_{\bullet}(\underline{F})$}
d'un Module $\underline{F}$ s'obtient en représentant $\underline{F}$ comme
quotient du Module engendré par le faisceau d'ensembles $\underline{F}$ :
$\varepsilon : L_0(\underline{F}) \to F$, en appliquant ensuite cette
construction à $\mathrm{Ker}(\varepsilon)$, et en itérant le procédé.
\marginal{réf ?}
\note{« $L_{\bullet}(\underline{F})$ » est ajouté à l'encre au-dessus de la
ligne ; « résolution plate canonique » est souligné au crayon, et un signe
d'insertion au crayon après « Rappelons » est relié à « réf ? ».}

Le foncteur $Rf^!$ est encore fourni, \add{par la méthode de 3.1.10,} par la
donnée
\[
U \longmapsto \add{K_0(U) \overset{\mathrm{df}}{=}}\; L_{\bullet}(f_!^{\bullet}(\mathbb{Z}/n)_U)\; \struck{= K_0(U)} .
\]
\note{les deux ajouts sont au crayon : « par la méthode de 3.1.10 » au-dessus
de la ligne, et « $K_0(U) \overset{\mathrm{df}}{=}$ » devant la formule, dont
la fin « $= K_0(U)$ », à l'encre, est biffée au crayon.}

Si maintenant $\underline{G}$ est un complexe de $\mathcal{A}$-Modules
acycliques ( ), on a
\[
K_0^{!}(\underline{G}) \xrightarrow{\ \sim\ } RK_0^{!}(\underline{G})
\]
\note{la parenthèse après « acycliques » est laissée en blanc ;
« acycliques, » est entouré au crayon, qui porte « ?? » dans la marge.}

On \struck{en déduit} \add{retrouve} que le faisceau d'anneaux $\mathcal{A}$
ne joue qu'un rôle bidon dans la construction de \struck{$RK_0^{!}$}
$Rf^!$, \struck{on}
\note{« retrouve » est à l'encre verte.}

\page{43}

\note{une bande de papier étroite, de deux lignes ; « cette construction »
renvoie vraisemblablement à celle de la page 42.}

Je n'ai pas vérifié que cette construction coïncide avec celle de 3.1.9.

\page{44}

3.1.15. On peut \struck{encore} décrire $Rf^!$ par le procédé de 3.1.13.
\struck{\ill{}} Soit en effet un morphisme compactifiable $f$ :

\begin{tikzcd}
X \arrow[r, hook, "j"'] \arrow[d, "f"'] & \bar X \arrow[dl, "\bar f"] \\
S &
\end{tikzcd}

et $\mathcal{A}$ un faisceau d'anneaux sur $S$ tel que $n\mathcal{A} = 0$
($n \geqslant 1$). Pour $W = (U, V, \varphi) \in \mathrm{Ob}\, X_f$
(notation de 3.1.13), soit
\[
(3.1.15.1) \qquad K(W) = \bar f_{V *}\, \tau_{\leqslant 2d}\, C^{*}_{\ell, \bar X_V}(\mathcal{A}_U)
\]
(résolution calculée sur $\bar X_V = \bar X \times_S V$).
\struck{On a un quasi-isomorphisme}
\note{le « 5 » de 3.1.15 est à l'encre verte sur un « 3 », et le « 5 » de
(3.1.15.1) repassé à l'encre noire ; le dernier chiffre de « 3.1.13 », dans
les deux renvois, est lui aussi repassé. Au crayon : $\tau_{\leqslant 2d}$ est
entouré et relié à « $\tau''$ ? » dans la marge de gauche, et
$\mathcal{A}_U$ est entouré et relié à « signification pas claire » dans la
marge de droite.}
\marginal{$\tau''$ ?}
\marginal{signification pas claire}

Pour $j$ \struck{inclusion} \add{morphisme structurel} de $V$ dans $S$, on a
un quasi-isomorphisme
\[
\struck{j_! K(W)}
\]
\[
j_!\, K(W) \xrightarrow{\ \sim\ } \bar f_{*}\, \tau_{\leqslant 2d}\, C^{*}_{\ell, \bar X}(\mathcal{A}_U)
\]
(XVII 5.2). D'après 3.1.9, le système (3.1.13) des $K(W)$ permet donc de
calculer $Rf^!$. \struck{Ceci donne un sens à la « formule d'induction »}
\note{« morphisme structurel » est ajouté à l'encre au-dessus de
« inclusion » biffé. Dans (3.1.13), les parenthèses et le dernier chiffre
sont à l'encre verte. Au crayon : $\bar f_{*}$ et $\tau_{\leqslant 2d}$ sont
entourés, et $\bar f_{*}$ relié à « $(j\bar f_V)_{*}$ ? » dans la marge ;
« ref » est porté devant (XVII 5.2), avec un signe d'insertion ; « 3.1.9 » et
« (3.1.13) » sont entourés, chacun avec un « ? », et reliés à « pas clair »
dans la marge de droite.}
\marginal{$(j\bar f_V)_{*}$ ?}
\marginal{ref}
\marginal{pas clair}

\note{dans la marge de gauche, au crayon, un diagramme des schémas en jeu. Les
deux traits verticaux qui joignent $X$ à $\bar X$ et $X_V$ à $\bar X_V$ n'ont
pas de pointe lisible ; les deux arcs portent leur pointe en bas.}

\begin{tikzcd}
\bar X \arrow[d, no head] \arrow[dd, bend right=50, "\bar f"'] & \bar X_V \arrow[d, no head] \arrow[dd, bend right=50, "\bar f_V"'] & \\
X \arrow[d, "f"'] & X_V \arrow[l] \arrow[d, "f_V"] & U \arrow[l] \\
S & V \arrow[l, "j"] &
\end{tikzcd}

\page{45}

\struck{Proposition 3.1.12. Le morphis}
\note{une seule ligne, en haut du feuillet ; le reste est blanc.}

\page{46}

3.1.16. L'isomorphisme \struck{d'induction} \add{d'adjonction} s'écrit
\[
(3.1.16.1) \qquad \mathrm{Hom}(Rf_!\, K, L) \simeq \mathrm{Hom}(K, Rf^! L)
\]
\add{avec} \struck{$\{$}$K \in \mathrm{Ob}\, D(X, f^{*}\mathcal{A})$ et
$L \in \mathrm{Ob}\, D^{+}(S, \mathcal{A})$.\struck{$\}$} Pour
$K \in \mathrm{Ob}\, D^{-}(X, f^{*}\mathcal{A})$ et
$L \in \mathrm{Ob}\, D^{+}(S, f^{*}\mathcal{A})$ \struck{et
\uncertain{isomorphe}} complexe d'injectifs, cet isomorphisme s'écrit, au
niveau des complexes,
\[
\mathrm{Hom}^{\bullet}(f_!^{\bullet} K, L) \simeq \mathrm{Hom}^{\bullet}(K, f^{!}_{\bullet} L) .
\]
\note{« d'adjonction » est écrit au-dessus de « d'induction » biffé ;
« avec » est à l'encre dans la marge. La seconde fois, la page écrit bien
$D^{+}(S, f^{*}\mathcal{A})$. Au crayon : « 3.1.16. » est entouré et relié à
la première note de marge ; « complexe d'injectifs » est entouré, et une
astérisque avant le mot biffé renvoie à la seconde ; la formule au niveau des
complexes est soulignée.}
\marginal{remarque \uncertain{en} cor. à 3.1.4 ?}
\marginal{$L$, ou $K$ et $L$ ?}

Pour tout \struck{$V \in \mathrm{Ob}\, S_{\text{et}}$}
\add{$j : V \to S$ dans $S_{\text{et}}$}, on dispose de
\[
\mathrm{Hom}^{*}_{\uncertain{/V}}(K, f^{!}_{\bullet} L)
= \mathrm{Hom}^{\bullet}(j'_! j'^{*} K, f^{!}_{\bullet} L)
\xrightarrow{\ \sim\ } \mathrm{Hom}^{\bullet}(f_!^{\bullet}(j'_! j'^{*} K), L) \longrightarrow
\]
\[
\struck{\longrightarrow \mathrm{Hom}\ \mathrm{Hom}}
\]
\[
\longrightarrow \mathrm{Hom}(j_! j^{*}(f_!^{\bullet} K), L) = \mathrm{Hom}_{V}(f_!^{\bullet} K, L)
\]
d'où un morphisme
\[
(3.1.16.2) \qquad Rf_{*}\, R\underline{\mathrm{Hom}}(K, Rf^! L) \xrightarrow{\ \sim\ } R\underline{\mathrm{Hom}}(Rf_!\, K, L)
\]
que par localisation de 3.1.16.1 on vérifie être un isomorphisme.
\note{l'indice du premier $\mathrm{Hom}$ et le petit signe écrit devant
l'indice $V$ du dernier ne se lisent pas sûrement. Un trait vertical au crayon
court en marge le long du calcul, avec la note qui suit.}
\marginal{devrait \uncertain{remonter} en cor. \uncertain{3.1.4.8} à 3.1.4
(conséquence triviale de \uncertain{3.1.4.7})}

3.1.17. \textbf{Exemple} : Soit $X$ un schéma compactifiable sur le spectre
d'un corps algébriquement clos $k$ \struck{\ill{}}, $x$ un point fermé de $X$
et $n$ un entier $\geqslant 1$ inversible dans $k$. La formule (3.1.11.1)
fournit, puisque $\mathbb{Z}/n\mathbb{Z}$ est un $\mathbb{Z}/n\mathbb{Z}$-module
injectif, \struck{Pour $U(x)$ parcourant les voisinages étales de $x$}
\[
\underline{H}^{-i}(Rf^! \mathbb{Z}/n\mathbb{Z})_x \simeq \varinjlim_{U(x)}
\mathrm{Hom}\bigl(\struck{\ill{}}\, H^{i}_{c}(U, \mathbb{Z}/n\mathbb{Z}), \mathbb{Z}/n\mathbb{Z}\bigr) ,
\]
\add{où $U(x)$ est la catégorie des voisinages étales de $x$ dans $X$.}
Il résulte vraisemblablement du théorème de bidualité (SGA 5 I )
(qui dépend de la résolution des singularités)
\note{le point après « dans $k$ » est à l'encre verte, et le « L » de « La »
au crayon ; le membre biffé « Pour $U(x)$ parcourant les voisinages étales de
$x$ » déborde dans la marge de gauche, et il est biffé à l'encre verte. La
ligne « où $U(x)$ est la catégorie… » est serrée entre deux lignes, d'un
trait plus pâle. Le numéro d'exposé après « SGA 5 I » est laissé en blanc.
Au crayon : « 3.1.17 » est entouré et relié à une note de marge, un « en » est
entouré au-dessus de « Soit », « vraisemblablement » est entouré, et un « ? »
est porté sous la parenthèse de la référence.}
\marginal{(\ill{}, on \ill{} en faisant \uncertain{venir} \ill{} d'une ligne)}
\marginal{\ill{} pas $W(x)$, $W$ ?}

\page{47}

que la flèche
\[
H^{i}_{x}(X, \mathbb{Z}/n\mathbb{Z}) \longrightarrow \struck{\ill{}}\;
\underset{U(x)}{\underset{\longleftarrow}{\text{“lim”}}}\; H^{i}_{c}(U, \mathbb{Z}/n\mathbb{Z})
\]
est un isomorphisme de pro-objets, le second membre étant représenté par le
premier. On aurait alors un isomorphisme
\[
\underline{H}^{-i}(Rf^! \mathbb{Z}/n\mathbb{Z})_x \xrightarrow{\ \sim\ } H^{i}_{x}(X, \mathbb{Z}/n\mathbb{Z})^{\vee}
\]
\note{un trait vertical au crayon borde en marge tout ce passage, depuis le
bas de la page 46.}

La propriété de « finitude » suivante résulte aussitôt \struck{Pro}
\add{de la construction} de \struck{ce que} des foncteurs $f^{!}_{i}$
(3.1.4.5) \add{et du fait qu'ils} transforment injectifs en injectifs.
\note{« de la construction » est ajouté à l'encre ; « et du fait qu'ils »
est au crayon.}

\textbf{Proposition 3.1.18.} (i) \struck{Si} \add{Supposons} le morphisme
compactifiable $f : X \to S$ est de dimension relative $\leqslant d$,
\struck{si \ill{}}, \struck{\ill{}} \add{et soit}
$L \in D^{+}(S, \mathcal{A})$\struck{, \ill{}} \add{tel que}
\struck{si} $\underline{H}^{i}(L) = 0$ pour $i \geqslant k$, alors
$\underline{H}^{i}(L) = 0$ pour $i \geqslant k - 2d$.

\struck{(ii) Si $\dim\mathrm{inj}(L) \leqslant k$, alors
$\dim\mathrm{inj}(Rf^! L) \leqslant k$\ill{}.}

(ii) $\dim\mathrm{inj}(Rf^! L) \leqslant \dim\mathrm{inj}(L)$.
\note{« Supposons », « et soit » et « tel que » sont au crayon, par-dessus ou
au-dessus des mots biffés. Dans la conclusion de (i), le second
$\underline{H}^{i}$ porte sur $L$, comme le premier : c'est ce que la page
écrit.}
\marginal{devrait \uncertain{remonter} avant 3.1.5 ! + immédiat \uncertain{après}}
\note{au crayon, entouré, en regard de la proposition ; un « \uncertain{ok} »
à l'encre rouge est écrit par-dessus, comme sur les deux notes de marge qui
suivent.}

3.1.19. Lorsque $f$ est quasi-fini et si on fait $d = 0$ dans la construction
3.1.4 du foncteur $Rf^!$, on trouve \add{par 3.1.3} que le foncteur $f_!$
admet un adjoint à droite $f^!$, dont $Rf^!$ est le dérivé. En particulier si
$f$ est une immersion $f : X \hookrightarrow S$, on trouve que $Rf^!$ est le
dérivé du foncteur « sections à support dans $X$ ». En particulier
\[
(3.1.19.1) \qquad R^{i} f^{!} L = \underline{H}^{i}_{X}(L)
\]
\note{« par 3.1.3 » est au crayon, entouré, avec « trouve que » ; un petit
signe au crayon est porté sous le $f$ de (3.1.19.1).}
\marginal{\uncertain{remonter} également avant 3.1.6. bref}
\marginal{(\ill{} dans \uncertain{qui} ?)}

\marginal{\uncertain{Dire} aussi que si $f$ est étale, $Rf^! = f^{*}$ ;
corollaire : « nature locale » de $Rf^!$}

\marginal{Le plan de 3.1 semble flou après 3.1.4. On pourrait prendre
a) Propriétés formelles de $Rf^!$ pour $f$ fixé : 3.1.5, 3.1.18, 3.1.19,
3.1.8, 3.1.9 ;
b) Propriétés formelles pour $f$ variable : 3.1.6 + [chgt de base par
morphisme lisse (\uncertain{univ.} loc. acyclique)] ;
c) Sorites sur construction locale de $Rf^!$ : 3.1.10 à 3.1.13.
(?)}
\note{tout le bas de la page est au crayon, sous un trait tiré en travers.
Le chiffre qui précède « 3.1.5 » dans a) est biffé. La mention entre crochets
de b) est encadrée, et une flèche y monte depuis « à ne pas oublier ! »,
souligné, écrit en bas à droite.}
\marginal{à ne pas oublier !}

\note{au pied de la page, toujours au crayon, un petit diagramme et un calcul
en une ligne. Le premier $\simeq$ du calcul est précédé d'un $f$ en trop ;
entre $\Delta^{*}(p_2^{!} f^{*})$ et la suite, quelques signes sont gommés ;
le $*$ de $f^{*}[2d](d)$ est écrit sur un « ! », et l'exposant du dernier $f$
ne se lit pas sûrement.}

\begin{tikzcd}
X \arrow[d, "f"'] & X \times_S X \arrow[l, "p_1"'] \arrow[d, "p_2"] & X \arrow[l, "\Delta"'] \\
S & X \arrow[l, no head, "f"] &
\end{tikzcd}

\[
f^{!} \simeq (p_1 \Delta)^{*} f^{!} \simeq \Delta^{*}(p_1^{*} f^{!})
\underset{\text{chgt de base par lisse}}{\simeq} \Delta^{*}(p_2^{!} f^{*})
\underset{\text{pureté relative ??}}{\simeq} \Delta^{!} p_2^{!} f^{*}[2d](d)
\simeq f^{\uncertain{*}}[2d](d)
\]

\end{document}
